\documentclass[11pt]{article}

\usepackage[margin=2.7cm]{geometry}
\usepackage{amsmath, amssymb}
\usepackage{booktabs}
\usepackage{enumitem}
\usepackage{xcolor}
\usepackage{hyperref}
\usepackage{parskip}
\usepackage{microtype}

\newcommand{\code}[1]{\texttt{#1}}
\newcommand{\ti}{t_i}
\newcommand{\tk}{t_k}
\newcommand{\va}{v_a}
\newcommand{\pd}[2]{\dfrac{\partial #1}{\partial #2}}

\hypersetup{colorlinks=true, linkcolor=blue!50!black, urlcolor=blue!50!black}

\title{Analytic Jacobian of the finite-difference SSA residual\\[2pt]
\large Companion to \code{momentum\_\allowbreak balance\_\allowbreak solver\_\allowbreak SSA\_\allowbreak FD\_\allowbreak SNES.f90}}
\author{}
\date{}

\begin{document}
\maketitle
\vspace{-2.5em}

\paragraph{Status.} Implemented and validated. Both the viscosity term
(Section~\ref{sec:visc}) and the sliding term (Section~\ref{sec:sliding}) are
assembled by \code{assemble\_\allowbreak SSA\_\allowbreak coeff\_\allowbreak jacobian\_\allowbreak CSR}, and added to the
Picard operator in \code{build\_\allowbreak SSA\_\allowbreak FD\_\allowbreak SNES\_\allowbreak jacobian\_\allowbreak petsc}. A central
finite-difference check of $Jv$ against $\big(F(u+hv)-F(u-hv)\big)/2h$ at a
non-zero base velocity gives a relative error of $\sim 8\times10^{-6}$
(Section~\ref{sec:validation}). Newton then converges quadratically on
MISMIP\_\allowbreak mod (residual history e.g.\ $12.5 \to 8.6 \to 4.8 \to 2.0 \to 0.57 \to
0.095 \to 0.0046$ in $\lVert \hat F\rVert$), where Tier~1/2 and the
50-iteration finite-difference Picard solver do not. The converged velocity
agrees with the Picard solver (run to 2000 iterations) to $\sim 5\%$ RMS, with
larger localised differences near the grounding line -- both solvers bottom
out at the noise floor of the non-smooth $\eta$ / friction / $\max(0.1,H)$
clamps. See \code{SSA\_\allowbreak FD\_\allowbreak SNES\_\allowbreak implementation\_\allowbreak plan.md} for the surrounding
context (Tiers~1--2, why they do not work, and the run log).

\tableofcontents

\section{Notation}
\label{sec:notation}

The SSA is discretised on two staggered entities of the mesh: the
\emph{a-grid} (mesh vertices) and the \emph{b-grid} (mesh triangles). Table
\ref{tab:notation} fixes the symbols used below and their corresponding
Fortran identifiers (all \code{type\_\allowbreak CSR\_\allowbreak matrix\_\allowbreak dp} sparse operators, or
plain fields, on \code{self\%mesh} / \code{self} unless noted).

\begin{table}[h]
\centering
\small
\begin{tabular}{@{}lp{5.0cm}p{5.9cm}@{}}
\toprule
Symbol & Fortran & Meaning \\
\midrule
$u,v$ & \code{u\_\allowbreak vav\_\allowbreak b}, \code{v\_\allowbreak vav\_\allowbreak b} & vertically-averaged velocity, b-grid (the unknowns) \\
$u_a,v_a$ & (local) & $u,v$ mapped to the a-grid \\
$\eta_a$ & \code{eta\_\allowbreak vav\_\allowbreak a} & effective viscosity, a-grid (after its min/max clamp) \\
$N_a = \bar H\,\eta_a$ & \code{N\_\allowbreak a} & viscosity--thickness product, a-grid \\
$N, N_x, N_y$ & \code{N\_\allowbreak b}, \code{dN\_\allowbreak dx\_\allowbreak b}, \code{dN\_\allowbreak dy\_\allowbreak b} & $N_a$ mapped to b-grid, and its gradient there \\
$\beta_b$ & \code{basal\_\allowbreak friction\_\allowbreak coefficient\_\allowbreak b} & applied basal friction coefficient, b-grid \\
$\tau_{d,x},\tau_{d,y}$ & \code{tau\_\allowbreak dx\_\allowbreak b}, \code{tau\_\allowbreak dy\_\allowbreak b} & driving stress, b-grid \\
$D_x, D_y$ & \code{M2\_\allowbreak ddx\_\allowbreak b\_\allowbreak b}, \code{M2\_\allowbreak ddy\_\allowbreak b\_\allowbreak b} & first-derivative operators, b$\to$b \\
$D_{xx}, D_{xy}, D_{yy}$ & \code{M2\_\allowbreak d2dx2\_\allowbreak b\_\allowbreak b}, \code{M2\_\allowbreak d2dxdy\_\allowbreak b\_\allowbreak b}, \code{M2\_\allowbreak d2dy2\_\allowbreak b\_\allowbreak b} & second-derivative operators, b$\to$b \\
$P_{a\to b}, P_{x,a\to b}, P_{y,a\to b}$ & \code{M\_\allowbreak map\_\allowbreak a\_\allowbreak b}, \code{M\_\allowbreak ddx\_\allowbreak a\_\allowbreak b}, \code{M\_\allowbreak ddy\_\allowbreak a\_\allowbreak b} & map / d/dx / d/dy, a$\to$b \\
$P_{b\to a}, P_{x,b\to a}, P_{y,b\to a}$ & \code{M\_\allowbreak map\_\allowbreak b\_\allowbreak a}, \code{M\_\allowbreak ddx\_\allowbreak b\_\allowbreak a}, \code{M\_\allowbreak ddy\_\allowbreak b\_\allowbreak a} & map / d/dx / d/dy, b$\to$a \\
\bottomrule
\end{tabular}
\caption{Symbol $\leftrightarrow$ code correspondence.}
\label{tab:notation}
\end{table}

A subscript operator applied to a field and evaluated at a mesh entity means
the matrix--vector product read at that row/entry, e.g.\
$u_x(\ti) := (D_x u)(\ti)$, and $P_{a\to b}(\ti,\va)$ is entry $(\ti,\va)$ of
that sparse operator (read via \code{\%read\_\allowbreak single\_\allowbreak row}). $\ti,\tk$ index
triangles (b-grid), $\va$ indexes vertices (a-grid).

\section{Newton's method and the Jacobian}
\label{sec:newton}

Let $n_u := 2\,n_{\mathrm{Tri}}$ be the number of unknowns (two velocity
components per triangle). Collect the unknowns on the b-grid into one vector
$u\in\mathbb{R}^{n_u}$, interleaved
$u=\big(u(t_1),v(t_1),u(t_2),v(t_2),\dots\big)$ in the \code{tiuv2n} ordering
the code uses throughout, and write the momentum-balance residual of
Section~\ref{sec:residual} as a map $F:\mathbb{R}^{n_u}\to\mathbb{R}^{n_u}$.
Solving the SSA means finding $u^\star$ with $F(u^\star)=0$.

Newton's method replaces this root-finding problem by a sequence of linear
ones. Given the current iterate $u_k$, a first-order Taylor expansion
linearises $F$ about it,
\begin{equation}
F(u_k+\delta u) \approx F(u_k) + J(u_k)\,\delta u,
\end{equation}
where $J(u)$, the \textbf{Jacobian}, is \emph{defined} as the matrix of partial
derivatives of every residual component with respect to every unknown,
\begin{equation}
J(u) := \frac{\partial F}{\partial u}, \qquad
J_{rc}(u) = \frac{\partial F_r(u)}{\partial u_c}, \qquad r,c = 1,\dots,n_u.
\label{eq:jacobian_def}
\end{equation}
Requiring the linearised residual to vanish, $F(u_k) + J(u_k)\,\delta u = 0$,
gives the \textbf{Newton step}
\begin{equation}
J(u_k)\,\delta u = -F(u_k), \qquad u_{k+1} = u_k + \delta u,
\label{eq:newton_step}
\end{equation}
which is exactly the linear system PETSc's SNES solves once per non-linear
iteration, from the values \code{SSA\_\allowbreak FD\_\allowbreak SNES\_\allowbreak form\_\allowbreak function} and
\code{SSA\_\allowbreak FD\_\allowbreak SNES\_\allowbreak form\_\allowbreak jacobian} return. \eqref{eq:jacobian_def} is a
generic definition -- it says nothing yet about the SSA specifically. The rest
of this document computes $J(u)$ for the residual of Section~\ref{sec:residual}
(Sections~\ref{sec:picard}--\ref{sec:sliding}), then checks the result
(Section~\ref{sec:validation}). In the code, $u$ and $F$ are further replaced
by the non-dimensional $\hat u = u/u_0$, $\hat F = F/F_0$ of
Section~\ref{sec:validation}; since $u_0,F_0$ are constants this only rescales
$J$ by $u_0/F_0$ and changes nothing derived below.

\section{The residual}
\label{sec:residual}

The SSA reads, in strong form,
\begin{align}
\frac{\partial}{\partial x}\Big[2N\big(2u_x+v_y\big)\Big] + \frac{\partial}{\partial y}\Big[N\big(u_y+v_x\big)\Big] - \beta_b\, u &= -\tau_{d,x}, \\
\frac{\partial}{\partial y}\Big[2N\big(2v_y+u_x\big)\Big] + \frac{\partial}{\partial x}\Big[N\big(v_x+u_y\big)\Big] - \beta_b\, v &= -\tau_{d,y}.
\end{align}
Expanded with the product rule and discretised with the b$\to$b operators of
Table~\ref{tab:notation} -- exactly as
\code{calc\_\allowbreak SSA\_\allowbreak DIVA\_\allowbreak stiffness\_\allowbreak matrix\_\allowbreak row\_\allowbreak free} does -- the residual on
triangle $\ti$ is
\begin{align}
F_x(\ti) &= 4N\,u_{xx} + 4N_x\,u_x + N\,u_{yy} + N_y\,u_y - \beta_b\,u
  + 3N\,v_{xy} + 2N_x\,v_y + N_y\,v_x + \tau_{d,x}, \label{eq:Fx}\\
F_y(\ti) &= 4N\,v_{yy} + 4N_y\,v_y + N\,v_{xx} + N_x\,v_x - \beta_b\,v
  + 3N\,u_{xy} + 2N_y\,u_x + N_x\,u_y + \tau_{d,y}, \label{eq:Fy}
\end{align}
where every term on the right is evaluated at $\ti$ (i.e.\ $N \equiv N(\ti)$,
$u_{xx}\equiv(D_{xx}u)(\ti)$, and so on), and $F_y$ is obtained from $F_x$ by
$x\leftrightarrow y$, $u\leftrightarrow v$.

\section{The Jacobian of \texorpdfstring{$F(u) = A(u)\,u - b(u)$}{F(u) = A(u) u - b(u)}}
\label{sec:picard}

Equations \eqref{eq:Fx}--\eqref{eq:Fy} are \emph{linear} in $u$ once the
coefficients $N,N_x,N_y,\beta_b$ (and the load $\tau_{d,x},\tau_{d,y}$) are
held fixed. Written as one matrix equation over all triangles,
\begin{equation}
F(u) = A(u)\,u - b(u),
\label{eq:F_matrix_form}
\end{equation}
where $A(u)$ is the velocity-\emph{dependent} coefficient matrix whose row
$\ti$ is exactly what \eqref{eq:Fx}--\eqref{eq:Fy} assembles, with
$N,N_x,N_y,\beta_b$ evaluated at the velocity $u$, and $b(u)$ collects the
(sign-flipped) driving stress $-\tau_{d,x},-\tau_{d,y}$ together with whatever
a boundary-condition row prescribes.

\subsection{Differentiating $A(u)\,u$: the product rule}

Apply the definition \eqref{eq:jacobian_def} to \eqref{eq:F_matrix_form}
entry-by-entry, for residual row $r$ and unknown column $c$:
\begin{align}
J_{rc}(u) = \pd{F_r(u)}{u_c}
  &= \pd{}{u_c}\bigg[\sum_{c'} A_{rc'}(u)\,u_{c'}\bigg] - \pd{b_r(u)}{u_c} \notag \\
  &= \sum_{c'} \pd{A_{rc'}(u)}{u_c}\,u_{c'} \;+\; \sum_{c'} A_{rc'}(u)\,\pd{u_{c'}}{u_c}
     \;-\; \pd{b_r(u)}{u_c}. \label{eq:product_rule}
\end{align}
The first term is the product rule applied to the \emph{coefficients}: how
does row $r$ change because $A$ itself depends on $u$. The second term is the
product rule applied to the \emph{unknown} $u_{c'}$ that $A$ multiplies; since
the unknowns are independent variables, $\partial u_{c'}/\partial u_c =
\delta_{c'c}$ (the Kronecker delta), so that sum collapses to the single term
$A_{rc}(u)$. In matrix form,
\begin{equation}
J(u) = A(u) + \Delta(u) - \pd{b(u)}{u}, \qquad
\Delta(u) := \sum_{c'} \pd{A_{\cdot,c'}(u)}{u}\,u_{c'},
\label{eq:jacobian_split}
\end{equation}
where $A_{\cdot,c'}$ denotes column $c'$ of $A$ (the $\cdot$ ranges over all
rows $r$), so $\partial A_{\cdot,c'}(u)/\partial u$ is that column
differentiated with respect to the whole vector $u$ -- a matrix with $(r,c)$
entry $\partial A_{rc'}(u)/\partial u_c$ -- and $\Delta(u)$'s own $(r,c)$ entry
is exactly the first term of \eqref{eq:product_rule}, $\Delta_{rc}(u) =
\sum_{c'} \big(\partial A_{rc'}(u)/\partial u_c\big)\,u_{c'}$.
$b(u)$ carries the driving stress, which does not depend on the velocity at
all, and on boundary rows either a prescribed constant or a value read from
the previous outer iterate\footnote{The \code{periodic\_\allowbreak ISMIP\_\allowbreak HOM} and
\code{infinite\_\allowbreak SSA\_\allowbreak icestream} boundary conditions reference
\code{u\_\allowbreak vav\_\allowbreak b\_\allowbreak prev}: within one Jacobian evaluation that is a fixed
snapshot, not a function of the $u$ this Jacobian is taken with respect to, so
the argument below still applies to those rows -- but see the
\emph{Idealised BCs} risk in \code{SSA\_\allowbreak FD\_\allowbreak SNES\_\allowbreak implementation\_\allowbreak plan.md}.}
-- in no case a function of the current unknown $u$. So $\partial b(u)/\partial u = 0$ and
\begin{equation}
\boxed{\,J(u) = A(u) + \Delta(u).\,}
\label{eq:jacobian_final}
\end{equation}

\subsection{Why $A(u)$ needs no new code}

The first term of \eqref{eq:jacobian_final} is, by \eqref{eq:jacobian_split},
literally the matrix obtained by evaluating $N,N_x,N_y,\beta_b$ at the current
velocity $u$ and \emph{freezing} them there -- which is exactly what a call to
\code{assemble\_\allowbreak SSA\_\allowbreak DIVA\_\allowbreak linearised\_\allowbreak matrix\_\allowbreak eq} /
\code{calc\_\allowbreak SSA\_\allowbreak DIVA\_\allowbreak stiffness\_\allowbreak matrix\_\allowbreak row\_\allowbreak free} already does, and
exactly what the hand-rolled viscosity (Picard) iteration already builds every
outer iteration to solve $A(u_k)\,u_{k+1}=b(u_k)$. This is not a coincidence
or a deliberate simplification: \eqref{eq:product_rule} shows it is a direct
consequence of differentiating a coefficient-matrix-times-unknown product --
one of the two product-rule terms is always the coefficient matrix itself,
because $\partial u_{c'}/\partial u_c=\delta_{c'c}$ collapses the sum over
columns back down to the matrix entry that was already there. So $A(u)$, the
\textbf{Picard operator}, is reused verbatim as one part of the true Newton
Jacobian, and only $\Delta(u)$ -- the coefficients' own sensitivity to $u$,
weighted by the current $u$ -- is new. Section~\ref{sec:visc} derives its
contribution from $N,N_x,N_y$ (the effective viscosity), Section~\ref{sec:sliding}
its contribution from $\beta_b$ (the sliding law); \code{assemble\_\allowbreak SSA\_\allowbreak coeff\_\allowbreak
jacobian\_\allowbreak CSR} assembles $\Delta(u)$ and
\code{build\_\allowbreak SSA\_\allowbreak FD\_\allowbreak SNES\_\allowbreak jacobian\_\allowbreak petsc} adds it to $A(u)$ per
\eqref{eq:jacobian_final}.

\section{Viscosity-derivative term}
\label{sec:visc}

$N$ depends on $u,v$ through the effective strain rate, so it contributes to
$\Delta(u)$ from \eqref{eq:jacobian_split}: the terms of
\eqref{eq:Fx}--\eqref{eq:Fy} that involve $N,N_x,N_y$ each pick up an extra
piece from differentiating those coefficients, on top of the Picard part of
Section~\ref{sec:picard}. Writing $w\in\{u,v\}$ for either unknown field,
\begin{align}
\pd{F_x(\ti)}{w(\tk)} \mathrel{+}= {}&
  \big[4u_{xx}+u_{yy}+3v_{xy}\big]\,\pd{N(\ti)}{w(\tk)}
  + \big[4u_x+2v_y\big]\,\pd{N_x(\ti)}{w(\tk)}
  + \big[u_y+v_x\big]\,\pd{N_y(\ti)}{w(\tk)}, \label{eq:visc_Fx}\\[4pt]
\pd{F_y(\ti)}{w(\tk)} \mathrel{+}= {}&
  \big[4v_{yy}+v_{xx}+3u_{xy}\big]\,\pd{N(\ti)}{w(\tk)}
  + \big[4v_y+2u_x\big]\,\pd{N_y(\ti)}{w(\tk)}
  + \big[v_x+u_y\big]\,\pd{N_x(\ti)}{w(\tk)}, \label{eq:visc_Fy}
\end{align}
with the bracketed factors -- the current velocity's derivatives at $\ti$ --
evaluated once per Jacobian assembly, and $\partial N/\partial w$ etc.\ common
to both equations, derived below.

\subsection{$\partial N/\partial w$ at a triangle}

$N,N_x,N_y$ on the b-grid are the a$\to$b image of $N_a$, so
\begin{align}
\pd{N(\ti)}{w(\tk)} &= \sum_{\va} P_{a\to b}(\ti,\va)\,\pd{N_a(\va)}{w(\tk)}, \label{eq:dN}\\
\pd{N_x(\ti)}{w(\tk)} &= \sum_{\va} P_{x,a\to b}(\ti,\va)\,\pd{N_a(\va)}{w(\tk)}, \\
\pd{N_y(\ti)}{w(\tk)} &= \sum_{\va} P_{y,a\to b}(\ti,\va)\,\pd{N_a(\va)}{w(\tk)}.
\end{align}

\subsection{$\partial N_a/\partial w$ at a vertex}

$N_a(\va) = \bar H(\va)\,\eta_a(\va)$ with $\bar H(\va) = \max\big(0.1,
H_i(\va)\big)$ (thickness-independent of $u,v$), and $\eta_a = \eta(E)$ with
$E$ the regularised effective strain-rate invariant,
\begin{align}
E(\va) &= u_x^2 + v_y^2 + u_x v_y + \tfrac14\big(u_y+v_x\big)^2 + \varepsilon_0, \\
\eta &= \tfrac12 A^{-1/n} E^{m}, \qquad m = \frac{1-n}{2n}, \\
\pd{\eta}{E} &= \frac{m\,\eta}{E},
\end{align}
with $n=$~\code{C\%Glens\_\allowbreak flow\_\allowbreak law\_\allowbreak exponent} and
$\varepsilon_0=$~\code{C\%Glens\_\allowbreak flow\_\allowbreak law\_\allowbreak epsilon\_\allowbreak sq\_\allowbreak 0} (fixed -- no
adaptive inflation, see \code{update\_\allowbreak SSA\_\allowbreak coefficients\_\allowbreak from\_\allowbreak velocity}), and
strain rates evaluated at $\va$ on the a-grid. The strain-rate sensitivities
of $E$ are
\begin{equation}
\pd{E}{u_x} = 2u_x+v_y, \qquad
\pd{E}{v_y} = 2v_y+u_x, \qquad
\pd{E}{u_y} = \pd{E}{v_x} = \tfrac12\big(u_y+v_x\big).
\end{equation}
The a-grid strain rates come from the b-grid velocity via
$u_x(\va) = (D_{x,b\to a}u)(\va)$, $u_y(\va)=(D_{y,b\to a}u)(\va)$, and
likewise for $v_x,v_y$, so
\begin{align}
\pd{N_a(\va)}{u(\tk)} &= \bar H(\va)\,\pd{\eta}{E}\bigg|_{\va}
  \Big[(2u_x+v_y)\big|_{\va}\,P_{x,b\to a}(\va,\tk) + \tfrac12(u_y+v_x)\big|_{\va}\,P_{y,b\to a}(\va,\tk)\Big], \\
\pd{N_a(\va)}{v(\tk)} &= \bar H(\va)\,\pd{\eta}{E}\bigg|_{\va}
  \Big[(2v_y+u_x)\big|_{\va}\,P_{y,b\to a}(\va,\tk) + \tfrac12(u_y+v_x)\big|_{\va}\,P_{x,b\to a}(\va,\tk)\Big].
\end{align}
Define, once per Jacobian assembly, the per-vertex scalars (on the a-grid,
gathered to all so that halo vertices in the $\sum_{\va} P_{a\to b}(\ti,\va)$
sums are available on every rank):
\begin{equation}
c(\va) = \bar H(\va)\,\pd{\eta}{E}\bigg|_{\va}
\quad\text{(0 where the $\eta$ clamp is active -- Section~\ref{sec:clamp})},
\end{equation}
\begin{equation}
g_{xx}(\va) = c(\va)\,(2u_x+v_y)\big|_{\va}, \qquad
g_{sh}(\va) = c(\va)\,\tfrac12(u_y+v_x)\big|_{\va}, \qquad
g_{yy}(\va) = c(\va)\,(2v_y+u_x)\big|_{\va},
\end{equation}
so that
\begin{align}
\pd{N_a(\va)}{u(\tk)} &= g_{xx}(\va)\,P_{x,b\to a}(\va,\tk) + g_{sh}(\va)\,P_{y,b\to a}(\va,\tk), \label{eq:dNa}\\
\pd{N_a(\va)}{v(\tk)} &= g_{yy}(\va)\,P_{y,b\to a}(\va,\tk) + g_{sh}(\va)\,P_{x,b\to a}(\va,\tk).
\end{align}

\subsection{Clamp consistency}
\label{sec:clamp}

\code{calc\_\allowbreak effective\_\allowbreak viscosity} clamps $\eta_a \leftarrow \min\big(\max(\eta_a,
\eta_{\min}), \eta_{\max}\big)$ with $\eta_{\min}=$~\code{C\%visc\_\allowbreak eff\_\allowbreak min}
and $\eta_{\max}$ recomputed as in that routine (with the fixed
$\varepsilon_0$). Where a bound is active, $\partial\eta/\partial E = 0$ by
construction of the clamp, so $c(\va)$ is set to $0$ when
$\eta_a(\va)\le\eta_{\min}$ or $\eta_a(\va)\ge\eta_{\max}$, keeping the
Jacobian consistent with the clamped $N$ used in the residual.

\section{Sliding-law term}
\label{sec:sliding}

The residual carries $-\beta_b(\ti)\,w(\ti)$ on the diagonal ($w=u$ for
\eqref{eq:Fx}, $w=v$ for \eqref{eq:Fy}), with $\beta_b$ frozen in the Picard
part $A(u)$. Since $\beta_b$ also depends on $u,v$, it contributes the second
piece of $\Delta(u)$ from \eqref{eq:jacobian_split}:
\begin{equation}
\pd{F_x(\ti)}{w(\tk)} \mathrel{+}= -u(\ti)\,\pd{\beta_b(\ti)}{w(\tk)}, \qquad
\pd{F_y(\ti)}{w(\tk)} \mathrel{+}= -v(\ti)\,\pd{\beta_b(\ti)}{w(\tk)}.
\label{eq:sliding_add}
\end{equation}

\subsection{Chain}

\code{calc\_\allowbreak applied\_\allowbreak basal\_\allowbreak friction\_\allowbreak coefficient} maps $u,v \to u_a,v_a$
($P_{b\to a}$), evaluates the sliding law on the a-grid to get $\beta_a$, maps
$\beta_a\to\beta_b$ ($P_{a\to b}$), then scales by the sub-grid grounded
fraction $\mathrm{fr}(\ti) = f_{\mathrm{gr},b}(\ti)^{\,p}$ (when
\code{C\%do\_\allowbreak GL\_\allowbreak subgrid\_\allowbreak friction}, with $p=$~\code{C\%subgrid\_\allowbreak friction\_\allowbreak }
\code{exponent\_\allowbreak on\_\allowbreak B\_\allowbreak grid}). Hence
\begin{equation}
\pd{\beta_b(\ti)}{w(\tk)} = \mathrm{fr}(\ti)\sum_{\va} P_{a\to b}(\ti,\va)\,\pd{\beta_a(\va)}{w(\tk)}.
\label{eq:dbeta_b}
\end{equation}

\subsection{Zoet--Iverson sliding law}

Only this law's derivative is implemented
(\code{calc\_\allowbreak sliding\_\allowbreak law\_\allowbreak ZoetIverson}); other laws make
\code{assemble\_\allowbreak SSA\_\allowbreak coeff\_\allowbreak jacobian\_\allowbreak CSR} \code{crash}. With
$q = 1/q_{\mathrm{ZI}}$, $u_t = u_{t,\mathrm{ZI}}$, $\delta_v$ the velocity
regularisation, and $\tau_c(\va) = $~\code{till\_\allowbreak yield\_\allowbreak stress}$(\va)$
(independent of $u,v$; $q_{\mathrm{ZI}}=$~\code{C\%slid\_\allowbreak ZI\_\allowbreak p},
$u_{t,\mathrm{ZI}}=$~\code{C\%slid\_\allowbreak ZI\_\allowbreak ut},
$\delta_v=$~\code{C\%slid\_\allowbreak delta\_\allowbreak v}),
\begin{align}
|u|(\va) &= \sqrt{\delta_v^2 + u_a(\va)^2 + v_a(\va)^2}, \\
\beta_a(\va) &= \tau_c(\va)\,|u|^{\,q-1}\,(|u|+u_t)^{-q}, \\
\pd{\beta_a}{|u|}\bigg|_{\va} &= \tau_c(\va)\,|u|^{\,q-2}\,(|u|+u_t)^{-q-1}\,
  \Big[(q-1)\,u_t - |u|\Big]. \label{eq:dbeta_dabs}
\end{align}
Since $u_a(\va) = (P_{b\to a}u)(\va)$,
\begin{align}
\pd{\beta_a(\va)}{u(\tk)} &= \pd{\beta_a}{|u|}\bigg|_{\va}\,\frac{u_a(\va)}{|u|(\va)}\,P_{b\to a}(\va,\tk), \label{eq:dbeta_a}\\
\pd{\beta_a(\va)}{v(\tk)} &= \pd{\beta_a}{|u|}\bigg|_{\va}\,\frac{v_a(\va)}{|u|(\va)}\,P_{b\to a}(\va,\tk).
\end{align}
Define, once per assembly on owned vertices then gathered,
\begin{equation}
s_u(\va) = \pd{\beta_a}{|u|}\bigg|_{\va}\frac{u_a(\va)}{|u|(\va)}, \qquad
s_v(\va) = \pd{\beta_a}{|u|}\bigg|_{\va}\frac{v_a(\va)}{|u|(\va)},
\end{equation}
so that \eqref{eq:dbeta_a} reads $\partial\beta_a(\va)/\partial u(\tk) =
s_u(\va)\,P_{b\to a}(\va,\tk)$, $\partial\beta_a(\va)/\partial v(\tk) =
s_v(\va)\,P_{b\to a}(\va,\tk)$.

\subsection{Why the sliding term matters}

Near $u=0$, $\partial\beta_a/\partial|u|$ in \eqref{eq:dbeta_dabs} is large and
velocity-weakening (its sign is fixed by $(q-1)u_t - |u| < 0$ for the
parameters used): omitting \eqref{eq:sliding_add} leaves Newton directions
badly wrong once the velocity develops away from zero, even though the
viscosity term alone is correct. With it, Newton is quadratic
(Section~\ref{sec:validation}). It is also why the solve is warm-started from
a few relaxed Picard iterations rather than from $u=0$: at exactly $u=0$ this
term is finite (as $u_a/|u|\to0$) but the sliding law is nearly
non-differentiable in its immediate neighbourhood, and Newton cannot get
started there.

\section{Validation}
\label{sec:validation}

Check $Jv$ against a central finite difference of the residual for a probe
direction $v$,
\begin{equation}
r = \frac{F(u+hv) - F(u-hv)}{2h},
\end{equation}
in the non-dimensional variables the SNES sees ($\hat u = u/u_0$, $\hat F =
F/F_0$ with $u_0=$~\code{velocity\_\allowbreak scale}, $F_0=$~\code{stress\_\allowbreak scale}), at a
\textbf{non-zero} base velocity. At $u=0$ both the viscosity
(Section~\ref{sec:visc}) and sliding (Section~\ref{sec:sliding}) derivatives
vanish identically while the finite difference still sees their onset over the
probe step, so the check is meaningless exactly at $u=0$. Away from it, with
$h=10^{-5}$,
\begin{equation}
\frac{\lVert Jv - r\rVert}{\lVert r\rVert} \approx 8\times10^{-6}.
\end{equation}

\section{Assembly notes (as implemented)}
\label{sec:assembly}

\begin{itemize}[leftmargin=1.4em]
\item \textbf{Stencil.} Row $\ti$ couples to $\tk$ via two hops ($\ti\to\va$
  through $P_{\cdot,a\to b}$, $\va\to\tk$ through $P_{\cdot,b\to a}$), a
  $\sim$2-ring of triangles -- wider than the Picard operator. \code{nnz\_\allowbreak est}
  is sized at \code{nTri\_\allowbreak loc * 2 * 250}.
\item \textbf{No duplicate CSR entries.} The coefficient-derivative CSR
  (\code{Jc\_\allowbreak CSR}, holding \eqref{eq:visc_Fx}--\eqref{eq:visc_Fy} together
  with \eqref{eq:sliding_add}) is built with a dense per-row accumulator
  (\code{rowbuf} + \code{hit} + \code{touch} list), emitting one
  \code{add\_\allowbreak entry} per touched column in ascending order (the
  \code{add\_\allowbreak entry} contract requires all of row $i$ before row $i+1$). It is
  converted to PETSc separately from the Picard operator and the two are
  combined with \code{MatAXPY}, so no CSR ever receives two entries for the
  same $(\text{row},\text{col})$.
\item \textbf{Distribution.} $P_{x,b\to a}$, $P_{y,b\to a}$, $P_{b\to a}$
  (\code{M\_\allowbreak ddx\_\allowbreak b\_\allowbreak a}, \code{M\_\allowbreak ddy\_\allowbreak b\_\allowbreak a}, \code{M\_\allowbreak map\_\allowbreak b\_\allowbreak a}) are gathered
  to full local copies once per solve (\code{gather\_\allowbreak CSR\_\allowbreak to\_\allowbreak all} =
  \code{gather\_\allowbreak to\_\allowbreak primary} + \code{MPI\_\allowbreak BCAST}), so
  \code{\%read\_\allowbreak single\_\allowbreak row} works for the halo vertices $\va$ that appear in
  the $\sum_{\va} P_{a\to b}(\ti,\va)$ sums but may not be owned by this rank.
  $u,v$ and the per-vertex $g_{xx},g_{sh},g_{yy},s_u,s_v$ factors are gathered
  as plain fields (mesh sizes here are small, so full copies per rank are
  cheap).
\item \textbf{Boundary / Dirichlet rows.} Rows where the Picard assembly
  writes a domain-border stencil or a Dirichlet identity get no
  coefficient-derivative term -- they are not the momentum residual.
\end{itemize}

\end{document}
